\documentclass[11pt]{article}

\usepackage[margin=1in]{geometry}
\usepackage{amsmath,amssymb,mathtools}
\usepackage{booktabs}

\setlength{\parindent}{0pt}
\setlength{\parskip}{0.55em}
\renewcommand{\arraystretch}{1.15}

\newcommand{\Oh}{\mathcal{O}}
\newcommand{\eoc}{\mathrm{EOC}}
\newcommand{\norm}[1]{\left\lVert #1\right\rVert}

\title{Local Truncation Error Handout\\
Finite Difference Methods for $u_t+a u_x=0$}
\author{Rishabh Kumar}
\date{\today}

\begin{document}
\maketitle

\section{Setup}

We consider
\begin{equation}
    u_t+a u_x=0,\qquad a>0,
\end{equation}
on the grid
\begin{equation}
    x_m=-2+mh,\qquad t_n=nk,\qquad
    v_m^n\approx u(x_m,t_n),
\end{equation}
with
\begin{equation}
    \lambda=\frac{k}{h},\qquad \nu=a\lambda=\frac{ak}{h}.
\end{equation}
For the assignment $a=1$, so $\nu=\lambda$.

\section{Discrete Error Norms}

At the final time $T=t_N$, define the grid error
\begin{equation}
    e_m^N=v_m^N-u(x_m,T).
\end{equation}
The numerical experiments use the following discrete norms:
\begin{align}
    E_1
    &=
    h\sum_m |e_m^N|,
    \\
    E_2
    &=
    \left(h\sum_m |e_m^N|^2\right)^{1/2},
    \\
    E_\infty
    &=
    \max_m |e_m^N|.
\end{align}
The first two are discrete analogues of the continuous $L^1$ and $L^2$ norms:
\begin{equation}
    \norm{e(\cdot,T)}_{L^1}
    =
    \int |e(x,T)|\,dx,
    \qquad
    \norm{e(\cdot,T)}_{L^2}
    =
    \left(\int |e(x,T)|^2\,dx\right)^{1/2}.
\end{equation}
The maximum norm $E_\infty$ measures the worst pointwise error.  It is therefore more
sensitive to localized oscillations or errors near nonsmooth corners.

The experimental order of convergence between two successive meshes $h$ and $h/2$ is
\begin{equation}
    \eoc
    =
    \frac{\log(E_h/E_{h/2})}{\log 2}.
\end{equation}
For the stability sweeps, we also use the final-time amplitude
\begin{equation}
    \norm{v^N}_\infty=\max_m |v_m^N|.
\end{equation}
Since the exact solution in the assignment satisfies $0\le u(x,t)\le1$, a computed
amplitude much larger than $1$ is clear evidence of numerical instability.

\section{Definition of Local Truncation Error}

The local truncation error is the residual obtained when the exact solution is
substituted into the finite difference scheme.  For example, if a scheme has the form
\begin{equation}
    L_{h,k}v_m^n=0,
\end{equation}
then the local truncation error is
\begin{equation}
    \tau_m^n=(L_{h,k}u)_m^n,
\end{equation}
where $u$ is the exact smooth solution of $u_t+a u_x=0$.  In this handout the schemes
are written in divided form, so $\tau_m^n$ has the same physical units as
$u_t+a u_x$.

If
\begin{equation}
    \tau_m^n=\Oh(k^p+h^q),
\end{equation}
then the method is locally consistent of order $p$ in time and $q$ in space.  For a
hyperbolic equation it is common to keep the Courant number fixed:
\begin{equation}
    k=\lambda h.
\end{equation}
Under this scaling,
\begin{equation}
    \Oh(k+h)=\Oh(h),\qquad
    \Oh(k+h^2)=\Oh(h),\qquad
    \Oh(k^2+h^2)=\Oh(h^2).
\end{equation}
This is why schemes with $\Oh(k+h)$ or $\Oh(k+h^2)$ local truncation error are
first-order methods in the usual refinement study, while schemes with
$\Oh(k^2+h^2)$ are second-order for smooth solutions.

The local truncation error is not the same as the final-time error.  The final error
also depends on stability, boundary closures, startup steps, and the regularity of the
exact solution.  In particular, the assigned initial data has a corner, so the measured
$E_\infty$ convergence rate can be lower than the formal smooth-solution order.

\section{Taylor Expansions}

The local truncation error calculations below use
\begin{align}
    u_m^{n+1}
    &=u+k u_t+\frac{k^2}{2}u_{tt}+\frac{k^3}{6}u_{ttt}+\Oh(k^4),\\
    u_m^{n-1}
    &=u-k u_t+\frac{k^2}{2}u_{tt}-\frac{k^3}{6}u_{ttt}+\Oh(k^4),\\
    u_{m+1}^n
    &=u+h u_x+\frac{h^2}{2}u_{xx}+\frac{h^3}{6}u_{xxx}
    +\frac{h^4}{24}u_{xxxx}+\Oh(h^5),\\
    u_{m-1}^n
    &=u-h u_x+\frac{h^2}{2}u_{xx}-\frac{h^3}{6}u_{xxx}
    +\frac{h^4}{24}u_{xxxx}+\Oh(h^5).
\end{align}
The PDE also implies
\begin{equation}
    u_t=-a u_x,\qquad u_{tt}=a^2u_{xx},\qquad u_{ttt}=-a^3u_{xxx}.
\end{equation}

\section{Explicit Schemes}

\subsection{FTFS}

\begin{equation}
    \frac{v_m^{n+1}-v_m^n}{k}
    +a\frac{v_{m+1}^n-v_m^n}{h}=0.
\end{equation}
Substitution gives
\begin{align}
    \tau_m^n
    &=
    \left(u_t+\frac{k}{2}u_{tt}+\Oh(k^2)\right)
    +a\left(u_x+\frac{h}{2}u_{xx}+\Oh(h^2)\right)\\
    &=
    \frac{k}{2}u_{tt}+\frac{ah}{2}u_{xx}
    +\Oh(k^2+h^2),
\end{align}
so
\begin{equation}
    \tau_m^n=\Oh(k+h).
\end{equation}

\subsection{FTBS}

\begin{equation}
    \frac{v_m^{n+1}-v_m^n}{k}
    +a\frac{v_m^n-v_{m-1}^n}{h}=0.
\end{equation}
Since
\begin{equation}
    \frac{u_m^n-u_{m-1}^n}{h}
    =
    u_x-\frac{h}{2}u_{xx}+\Oh(h^2),
\end{equation}
we obtain
\begin{equation}
    \tau_m^n
    =
    \frac{k}{2}u_{tt}
    -\frac{ah}{2}u_{xx}
    +\Oh(k^2+h^2)
    =
    \Oh(k+h).
\end{equation}

\subsection{FTCS}

\begin{equation}
    \frac{v_m^{n+1}-v_m^n}{k}
    +a\frac{v_{m+1}^n-v_{m-1}^n}{2h}=0.
\end{equation}
The centered derivative satisfies
\begin{equation}
    \frac{u_{m+1}^n-u_{m-1}^n}{2h}
    =
    u_x+\frac{h^2}{6}u_{xxx}+\Oh(h^4).
\end{equation}
Therefore
\begin{equation}
    \tau_m^n
    =
    \frac{k}{2}u_{tt}
    +\frac{a h^2}{6}u_{xxx}
    +\Oh(k^2+h^4)
    =
    \Oh(k+h^2).
\end{equation}
With $k=\Oh(h)$, FTCS is first-order consistent.

\subsection{CTCS or Leapfrog}

\begin{equation}
    \frac{v_m^{n+1}-v_m^{n-1}}{2k}
    +a\frac{v_{m+1}^n-v_{m-1}^n}{2h}=0.
\end{equation}
The centered time difference gives
\begin{equation}
    \frac{u_m^{n+1}-u_m^{n-1}}{2k}
    =
    u_t+\frac{k^2}{6}u_{ttt}+\Oh(k^4).
\end{equation}
Hence
\begin{equation}
    \tau_m^n
    =
    \frac{k^2}{6}u_{ttt}
    +\frac{a h^2}{6}u_{xxx}
    +\Oh(k^4+h^4)
    =
    \Oh(k^2+h^2).
\end{equation}

\subsection{Lax--Friedrichs}

\begin{equation}
    \frac{v_m^{n+1}-\frac12(v_{m+1}^n+v_{m-1}^n)}{k}
    +a\frac{v_{m+1}^n-v_{m-1}^n}{2h}=0.
\end{equation}
Since
\begin{equation}
    \frac12(u_{m+1}^n+u_{m-1}^n)
    =
    u+\frac{h^2}{2}u_{xx}+\Oh(h^4),
\end{equation}
we obtain
\begin{equation}
    \tau_m^n
    =
    \frac{k}{2}u_{tt}
    -\frac{h^2}{2k}u_{xx}
    +\frac{a h^2}{6}u_{xxx}
    +\Oh(k^2+h^4/k+h^4).
\end{equation}
Thus
\begin{equation}
    \tau_m^n=\Oh\left(k+\frac{h^2}{k}\right).
\end{equation}
For $k=\Oh(h)$ this becomes first order.

\subsection{Lax--Wendroff}

\begin{equation}
    v_m^{n+1}
    =
    v_m^n
    -\frac{\nu}{2}(v_{m+1}^n-v_{m-1}^n)
    +\frac{\nu^2}{2}(v_{m+1}^n-2v_m^n+v_{m-1}^n).
\end{equation}
Equivalently,
\begin{equation}
    \frac{v_m^{n+1}-v_m^n}{k}
    +aD_0v_m^n-\frac{a^2k}{2}\delta_{xx}v_m^n=0.
\end{equation}
Using
\begin{equation}
    D_0u=u_x+\frac{h^2}{6}u_{xxx}+\Oh(h^4),
    \qquad
    \delta_{xx}u=u_{xx}+\frac{h^2}{12}u_{xxxx}+\Oh(h^4),
\end{equation}
we get
\begin{align}
    \tau_m^n
    &=
    u_t+a u_x
    +\frac{k}{2}u_{tt}
    -\frac{a^2k}{2}u_{xx}
    +\frac{k^2}{6}u_{ttt}
    +\frac{a h^2}{6}u_{xxx}
    +\Oh(k^3+kh^2+h^4).
\end{align}
The PDE cancels $u_t+a u_x$, and $u_{tt}=a^2u_{xx}$ cancels the order-$k$ term.  Therefore
\begin{equation}
    \tau_m^n=\Oh(k^2+h^2).
\end{equation}

\subsection{Beam--Warming}

For $a>0$,
\begin{equation}
    v_m^{n+1}
    =
    v_m^n
    -\frac{\nu}{2}(3v_m^n-4v_{m-1}^n+v_{m-2}^n)
    +\frac{\nu^2}{2}(v_m^n-2v_{m-1}^n+v_{m-2}^n).
\end{equation}
The backward approximations are
\begin{align}
    \frac{3u_m-4u_{m-1}+u_{m-2}}{2h}
    &=
    u_x-\frac{h^2}{3}u_{xxx}+\Oh(h^3),\\
    \frac{u_m-2u_{m-1}+u_{m-2}}{h^2}
    &=
    u_{xx}-h u_{xxx}+\Oh(h^2).
\end{align}
Substitution gives
\begin{equation}
    \tau_m^n
    =
    u_t+a u_x
    +\frac{k}{2}u_{tt}
    -\frac{a^2k}{2}u_{xx}
    +\Oh(k^2+h^2+kh).
\end{equation}
The first two nonzero order-$k$ terms cancel by $u_{tt}=a^2u_{xx}$, so
\begin{equation}
    \tau_m^n=\Oh(k^2+h^2)
\end{equation}
when $k=\Oh(h)$.

\section{Implicit Schemes}

\subsection{BTFS and BTBS}

BTFS and BTBS are respectively
\begin{align}
    \frac{v_m^{n+1}-v_m^n}{k}
    +a\frac{v_{m+1}^{n+1}-v_m^{n+1}}{h}
    &=0,\\
    \frac{v_m^{n+1}-v_m^n}{k}
    +a\frac{v_m^{n+1}-v_{m-1}^{n+1}}{h}
    &=0.
\end{align}
Expanding around $(x_m,t_{n+1})$,
\begin{equation}
    \frac{u_m^{n+1}-u_m^n}{k}
    =
    u_t-\frac{k}{2}u_{tt}+\Oh(k^2),
\end{equation}
while the forward/backward spatial differences are first-order accurate.  Thus
\begin{equation}
    \tau_m^{n+1}=\Oh(k+h).
\end{equation}

\subsection{BTCS}

\begin{equation}
    \frac{v_m^{n+1}-v_m^n}{k}
    +a\frac{v_{m+1}^{n+1}-v_{m-1}^{n+1}}{2h}=0.
\end{equation}
The time difference is first-order and the spatial difference is second-order:
\begin{equation}
    \tau_m^{n+1}
    =
    -\frac{k}{2}u_{tt}
    +\frac{a h^2}{6}u_{xxx}
    +\Oh(k^2+h^4)
    =
    \Oh(k+h^2).
\end{equation}

\subsection{Crank--Nicolson}

\begin{equation}
    \frac{v_m^{n+1}-v_m^n}{k}
    +
    \frac{a}{2}
    \left[
        \frac{v_{m+1}^{n+1}-v_{m-1}^{n+1}}{2h}
        +
        \frac{v_{m+1}^{n}-v_{m-1}^{n}}{2h}
    \right]=0.
\end{equation}
This is trapezoidal in time and centered in space.  Expanding about the midpoint $t_{n+1/2}$ gives
\begin{equation}
    \tau_m^{n+1/2}=\Oh(k^2+h^2).
\end{equation}

\section{Summary}

\begin{table}[h]
\centering
\begin{tabular}{@{}lll@{}}
\toprule
Scheme & Local truncation error & Order for $k=\Oh(h)$\\
\midrule
FTFS & $\Oh(k+h)$ & first\\
FTBS & $\Oh(k+h)$ & first\\
FTCS & $\Oh(k+h^2)$ & first\\
CTCS & $\Oh(k^2+h^2)$ & second\\
Lax--Friedrichs & $\Oh(k+h^2/k)$ & first\\
Lax--Wendroff & $\Oh(k^2+h^2)$ & second\\
Beam--Warming & $\Oh(k^2+h^2)$ & second\\
BTFS & $\Oh(k+h)$ & first\\
BTBS & $\Oh(k+h)$ & first\\
BTCS & $\Oh(k+h^2)$ & first\\
Crank--Nicolson & $\Oh(k^2+h^2)$ & second\\
\bottomrule
\end{tabular}
\end{table}

\end{document}
